% status: FULL
\chapter{Reproducing the Book's Measurements}\label{app:reproducing}
\begin{ChapterIntent}
Reproduction executes the stated experiment and compares observations;
it does not demand identical speed from different hardware. This chapter
routes readers to canonical evidence rather than copying raw logs.
The foundational arithmetic runs without a model download.
\end{ChapterIntent}

\section{Begin with a question that can be wrong}
Suppose you believe caching should make the decoder faster. That
sentence leaves several questions open. Faster for which prompt
length? Does the timer include setup? Are the cached and uncached
paths producing equivalent logits? How many tokens are generated?
An experiment becomes useful when those questions have answers.

For the Part I decoder, first test the semantic claim: at each
prefix, cached and full-prefix execution produce matching logits.
Only then ask the performance question. This ordering prevents an
incorrect shortcut from winning a benchmark by doing less than the
required computation.

A good record states the prediction before the result. For example:
as the prefix grows, reusing earlier layer state should avoid
increasing amounts of recomputation. That prediction can fail to
produce a practical speedup if bookkeeping overhead dominates at
tiny sizes. Keeping the negative result teaches something about
the implementation and workload.

\section{State the hypothesis}
Batching reuses weights is a mechanism claim. A batch improves goodput
at a latency target is a workload claim requiring arrivals, lengths and
an acceptance rule. An experiment should be able to falsify the latter,
not merely print an attractive throughput.

Measured numbers name hardware, software revision, model bytes, precision,
workload, warm-up, repetitions and statistic. Derived numbers name inputs
and assumptions. Illustrative numbers teach without claiming a device
result. A vendor peak stays a specification when inserted into a formula.

\section{An auditable timestamp trace}
Consider an illustrative request with times in milliseconds from intended
arrival: client send at 3, admission at 7, first delivered token at 17,
later tokens at 25 and 40, terminal event at 44.
TTFT from send is $17-3=14$ ms; from intended arrival it is 17 ms.
The first definition excludes 3 ms of client queueing.

Inter-token gaps are 8 and 15 ms. Mean time per output token excluding
the first is $(40-17)/(3-1)=11.5$ ms.
Terminal latency from send is $44-3=41$ ms. Using terminal time in the
token-gap formula would charge terminal handling to decode. Neither clock
is wrong; an unnamed boundary is wrong. A one-token response has no
inter-token sample under this convention.

The percentile of pooled gaps need not equal the percentile of per-request
mean gaps. State the population. Excluding failed or timed-out requests
can improve the reported tail while the service deteriorates.
Report errors and useful completed output alongside latency.

\section{Draw the clock before calculating latency}
The illustrative trace earlier names seven different events. Put them
on one line before reducing them to a statistic.
\EngineFigure{foundation-measurement-clock}{Illustrative client and
request timestamps in milliseconds. TTFT from send is 14 ms, while
the two delivered-token gaps are 8 and 15 ms. The terminal event
occurs after the last token and is not another token arrival.}
{fig:foundation-measurement-clock}

The intended arrival at zero expresses the workload schedule.
The actual send at three reveals a client-side delay. Admission at
seven describes a server boundary. The first token at seventeen
describes an observed delivery. These timestamps should not be
substituted for one another because they happen to be close in a
particular run.

Run the trace arithmetic:
\begin{verbatim}
python3 code/foundations/lesson.py measurement
\end{verbatim}
The program reports TTFT 14 ms, gaps 8 and 15 ms, mean gap 11.5 ms
and terminal latency 41 ms. Now deliberately use the terminal event
as the last token: $(44-17)/2=13.5$ ms. The arithmetic is valid,
but it answers a different question and mislabels the result if
reported as mean delivered-token gap.

Avoid subtracting timestamps from unsynchronized machines as though
they share one clock. End-to-end duration can be measured on the
client's monotonic clock. Internal stage durations can be measured
within their own clock domain. Combining them requires an explicit
clock and synchronization method.

\section{One run is an observation, not a distribution}
Repeated runs reveal variation from scheduling, caches, temperature,
background work and the workload itself. Report a central value and
a measure of spread appropriate to the question. Keep the raw
observations so another reader can recompute the summary.

For the illustrative times $[9,10,10,11,60]$ ms, the median is 10 ms
and the mean is 20 ms. The slow run is part of the data. Removing it
because it makes the chart unattractive changes the claim. If an
external interruption invalidated a run, state the exclusion rule
and retain the excluded observation with its reason.

Do not infer a reliable extreme-tail percentile from a handful of
samples. Percentile conventions also differ in interpolation and
rounding. State the method when exact reproduction matters. For a
service, retain failures and timeouts in the workload accounting;
otherwise the system can appear faster by failing its slowest requests.

\section{Warm means something specific}
A process can be new while the operating system still holds its
checkpoint pages in memory. A model can be loaded while its kernels
are not yet compiled. A kernel can be compiled while a particular
graph shape has not been captured. These are different warm states.

Describe which work has already happened before starting the timer.
For an isolated kernel experiment, preallocated buffers and compiled
kernels may be appropriate. For startup latency, excluding them
would remove the very cost being studied.

Do not flush shared caches, change system protections or download
large checkpoints merely to make an experiment look more controlled.
Use the available environment, document its limits, and separate
the question you can answer from the one that needs another setup.
A reproducible limitation is better evidence than an undocumented
intervention.

\section{A load generator can change the workload}
In a closed-loop test, a client sends new work only after earlier work
finishes. A slow service therefore receives fewer new requests.
In an open-loop test, arrivals follow an external schedule, so a
slowdown can build a queue. Both are useful, but they model different
users and produce different latency distributions.

The load generator can also fall behind. Record intended arrival,
actual send and observed completion so that client delay is visible.
If the client stops creating work during a server stall, measuring
only sent requests can hide the waiting that a real arrival stream
would experience.

For the same reason, report both throughput and useful completed
work under a latency target. A larger batch can raise aggregate
tokens per second while making each user wait longer. The systems
chapters will call useful work meeting the target \emph{goodput}.
The definition includes the acceptance rule; it is not another name
for the largest throughput number on a chart.

\section{Control arrivals and machine state}
A closed-loop client sends its next request after a response. An open-loop
client schedules arrivals independently. During a stall the first reduces
load, while the second can build a queue. Capture intended arrival and
actual send if the client itself can bottleneck. Quietly ceasing arrivals
during overload tests a different workload.

Compilation, graph capture, caches, thermal limits and competing work
change results. Separate warm-up, repeat runs and report spread as well
as a central statistic. Do not remove outliers without a stated rule and
reason. A cold process does not establish a cold OS file cache.

Preserve negative outcomes: a slower optimization or failed oracle is
evidence. Reproduction and layout inspection do not promote a chapter to
VERIFIED; that additionally requires independent technical review.

\section{Leave a record another reader can follow}
A compact record should let someone reconstruct the experiment
without guessing. Include the code revision, model identity,
hardware and software versions, command, input construction,
timer boundaries, warm-up, repetitions, summary method and raw
outputs. Separate facts you measured from specifications and
derived estimates used to interpret them.

When a command has several modes, record the arguments actually
used. A source file containing a GPU variant does not prove the
reported run used it. Similarly, a passing unit test does not prove
a production request reaches the tested code path.

Preserve the result when an optimization loses. Explain the initial
hypothesis, the changed variable and the observed outcome. That
record can be more valuable than a successful benchmark because it
shows where an apparently sensible cost model was incomplete.

\section{One home for each kind of evidence}
\begin{longtable}{@{}p{.28\linewidth}p{.63\linewidth}@{}}
\toprule Material & Canonical location and purpose\\\midrule
\endhead
Systems deepening & \texttt{research/measurements/2026-09-26-m1-pass3.md};
sections use stable source IDs, not the new printed numbering\\
Part experiments & Dated \texttt{m1-part2} through \texttt{m1-part8} records
in \texttt{research/measurements}; retain their limitations\\
Model byte accounting & \texttt{2026-09-23-gguf-tensor-bytes.md}
in the same directory\\
Original construction & \texttt{research/benchmarks} and \texttt{labs};
first-edition lab numbers remain unchanged\\
Current small labs & \texttt{code/labs/chNN-*}; README instructions for
prediction, execution, explanation and deliberate breakage\\
Primary-source facts & \texttt{research/FRONTIER.md}; source, check date
and implementation maturity\\
\bottomrule
\end{longtable}

The M1 records cover CPU and Metal experiments, not independent measurements
of every accelerator mentioned. Read each record's load and hardware notes.
An industrial engine's published result is not a locally reproduced result.
A progress transcript is not a substitute for a committed measurement
record. Lab directory numbers are stable identifiers, not printed chapter
numbers in the restructured edition.

\Needspace{10\baselineskip}
\section{Run cheap checks before expensive measurements}
From the repository root:
\begin{verbatim}
python3 scripts/check-foundations.py
make textbook-check
(cd code/mini-engine && cargo test --workspace)
\end{verbatim}
These check the worked examples, book contracts and teaching engine.
They do not establish benchmark performance. For a chapter experiment,
read its README and record, inspect the arguments, run the oracle, then
time the workload. Keep raw stdout and stderr including failures.
Record model content hashes and source revision, not only mutable aliases.

Do not download additional checkpoints, purge shared caches or disable
system protections as an unannounced part of reproduction. State any
unavailable hardware or unperformed experiment explicitly.

\section{From building a model to explaining an engine}
Part I has now connected the model's computation to its numerical
representation, hardware execution and evidence. You can follow
a token through the decoder, count its state, read its stored
weights and design a test of a claimed improvement.

The next part asks why a correct decoder is not yet an efficient
inference engine. Its measurements will be easier to read now:
every byte, operation and timestamp belongs to an object already
introduced. Keep the small model and its tests nearby. They remain
the reference when the implementation becomes more sophisticated.
\section{Worked problems}
\subsection*{1. Name the endpoints}
\textbf{Problem.} Send is at 3 ms, first token at 17 ms, last token at
40 ms and terminal at 44 ms. There are three tokens. Find TTFT and mean gap.
\textbf{Solution.} TTFT is 14 ms. Mean delivered-token gap is
$(40-17)/(3-1)=11.5$ ms. The terminal event is excluded from that mean.
\subsection*{2. Summarize without deleting}
\textbf{Problem.} Find the mean and median of $[9,10,10,11,60]$ ms.
\textbf{Solution.} The sum is 100 ms, so the mean is 20 ms.
The middle observation is 10 ms. Retain the 60 ms observation unless
a declared exclusion rule invalidates it.
\subsection*{3. A fresh process}
\textbf{Problem.} Does restarting the model process prove a cold file cache?
\textbf{Solution.} No. File pages can remain in the operating system's
cache after the process exits. State process and cache conditions separately.

